% GENERATED FILE. Edit canonical lessons, then run npm run generate:books.
% canonical-source-hash: fnv1a64:71dc525c3e140abd
% canonical-lessons: MR-C186-papad, MR-C186-catni, MR-C186-kosimbir, MR-C186-sira, MR-C186-modak

\chapter{Papad, Chutney, Salad, Semolina pudding, Modak}
\label{ch:mr-papad-chutney-salad-semolina-pudding-modak}
\hlchaptermodality{186}

\begin{chapteropening}
\textbf{By the end of this chapter:} \emph{I can say a papad, a chutney, a salad, a semolina pudding and a modak.}

Complete the last lesson of chapter 186: I can say a papad, a chutney, a salad, a semolina pudding and a modak.
\end{chapteropening}

\section[\texorpdfstring{\mr{पापड} (pāpaḍ)}{पापड (pāpaḍ)}]{\mr{पापड} (pāpaḍ) — a papad}

\label{lesson:MR-C186-papad}



\begin{quote}
Before the new one: say the Marathi for a defeat, then the Marathi for a competition.
\end{quote}

\subsection*{You'll want to know: \mr{पापड}}
\textbf{\mr{पापड}} — \emph{pāpaḍ} — \textquotedblleft{}a papad\textquotedblright{}.

\begin{grammarlens}[title={What you've built}]
One more word for this chapter.
\end{grammarlens}

\subsection*{Your turn}
\begin{itemize}
\raggedright
  \item \emph{Say it:} \emph{pāpaḍ}
  \item \emph{Say it:} \emph{pāpaḍ}, once more
  \item \emph{Say it:} say \emph{pāpaḍ}
  \item \emph{Recall:} say the Marathi for a defeat, then the Marathi for a competition, then say \emph{pāpaḍ} again
\end{itemize}

\begin{tcolorbox}[breakable,colback=teal!4,colframe=teal!35!black,title={Before you move on}]
What does \mr{पापड} mean? (\textquotedblleft{}A papad\textquotedblright{}.) Say it once more.
\end{tcolorbox}
\section[\texorpdfstring{\mr{चटणी} (caṭṇī)}{चटणी (caṭṇī)}]{\mr{चटणी} (caṭṇī) — a chutney}

\label{lesson:MR-C186-catni}



\begin{quote}
Before the new one: say the Marathi for a competition, then the Marathi for a papad.
\end{quote}

\subsection*{You'll want to know: \mr{चटणी}}
\textbf{\mr{चटणी}} — \emph{caṭṇī} — \textquotedblleft{}a chutney\textquotedblright{}.

\begin{grammarlens}[title={What you've built}]
One more word for this chapter.
\end{grammarlens}

\subsection*{Your turn}
\begin{itemize}
\raggedright
  \item \emph{Say it:} \emph{caṭṇī}
  \item \emph{Say it:} \emph{caṭṇī}, once more
  \item \emph{Say it:} say \emph{caṭṇī}
  \item \emph{Recall:} say the Marathi for a competition, then the Marathi for a papad, then say \emph{caṭṇī} again
\end{itemize}

\begin{tcolorbox}[breakable,colback=teal!4,colframe=teal!35!black,title={Before you move on}]
What does \mr{चटणी} mean? (\textquotedblleft{}A chutney\textquotedblright{}.) Say it once more.
\end{tcolorbox}
\section[\texorpdfstring{\mr{कोशिंबीर} (kośimbīr)}{कोशिंबीर (kośimbīr)}]{\mr{कोशिंबीर} (kośimbīr) — a salad}

\label{lesson:MR-C186-kosimbir}



\begin{quote}
Before the new one: say the Marathi for a papad, then the Marathi for a chutney.
\end{quote}

\subsection*{You'll want to know: \mr{कोशिंबीर}}
\textbf{\mr{कोशिंबीर}} — \emph{kośimbīr} — \textquotedblleft{}a salad\textquotedblright{}.

\begin{grammarlens}[title={What you've built}]
One more word for this chapter.
\end{grammarlens}

\subsection*{Your turn}
\begin{itemize}
\raggedright
  \item \emph{Say it:} \emph{kośimbīr}
  \item \emph{Say it:} \emph{kośimbīr}, once more
  \item \emph{Say it:} say \emph{kośimbīr}
  \item \emph{Recall:} say the Marathi for a papad, then the Marathi for a chutney, then say \emph{kośimbīr} again
\end{itemize}

\begin{tcolorbox}[breakable,colback=teal!4,colframe=teal!35!black,title={Before you move on}]
What does \mr{कोशिंबीर} mean? (\textquotedblleft{}A salad\textquotedblright{}.) Say it once more.
\end{tcolorbox}
\section[\texorpdfstring{\mr{शिरा} (śirā)}{शिरा (śirā)}]{\mr{शिरा} (śirā) — a semolina pudding}

\label{lesson:MR-C186-sira}



\begin{quote}
Before the new one: say the Marathi for a chutney, then the Marathi for a salad.
\end{quote}

\subsection*{You'll want to know: \mr{शिरा}}
\textbf{\mr{शिरा}} — \emph{śirā} — \textquotedblleft{}a semolina pudding\textquotedblright{}.

\begin{grammarlens}[title={What you've built}]
One more word for this chapter.
\end{grammarlens}

\subsection*{Your turn}
\begin{itemize}
\raggedright
  \item \emph{Say it:} \emph{śirā}
  \item \emph{Say it:} \emph{śirā}, once more
  \item \emph{Say it:} say \emph{śirā}
  \item \emph{Recall:} say the Marathi for a chutney, then the Marathi for a salad, then say \emph{śirā} again
\end{itemize}

\begin{tcolorbox}[breakable,colback=teal!4,colframe=teal!35!black,title={Before you move on}]
What does \mr{शिरा} mean? (\textquotedblleft{}A semolina pudding\textquotedblright{}.) Say it once more.
\end{tcolorbox}
\section[\texorpdfstring{\mr{मोदक} (modak)}{मोदक (modak)}]{\mr{मोदक} (modak) — a modak}

\label{lesson:MR-C186-modak}



\begin{quote}
Before the new one: say the Marathi for a salad, then the Marathi for a semolina pudding.
\end{quote}

\subsection*{You'll want to know: \mr{मोदक}}
\textbf{\mr{मोदक}} — \emph{modak} — \textquotedblleft{}a modak\textquotedblright{}.

\begin{grammarlens}[title={What you've built}]
One more word for this chapter.
\end{grammarlens}

\subsection*{Your turn}
\begin{itemize}
\raggedright
  \item \emph{Say it:} \emph{modak}
  \item \emph{Say it:} \emph{modak}, once more
  \item \emph{Say it:} say \emph{modak}
  \item \emph{Recall:} say the Marathi for a salad, then the Marathi for a semolina pudding, then say \emph{modak} again
\end{itemize}

\begin{tcolorbox}[breakable,colback=teal!4,colframe=teal!35!black,title={Before you move on}]
What does \mr{मोदक} mean? (\textquotedblleft{}A modak\textquotedblright{}.) Say it once more.
\end{tcolorbox}
